% Solutions of 23-dynamic-programming.
\section*{Chapter \ref{ch:b1:dynamic-programming} --- Dynamic Programming}

\begin{solution}{exo:b1:dynamic-programming:1}
Rows for the prefixes of \code{loan}, columns for those of \code{lane}:
\begin{center}
\begin{tabular}{c|ccccc}
 & & l & a & n & e \\ \hline
 & 0 & 1 & 2 & 3 & 4 \\
l & 1 & 0 & 1 & 2 & 3 \\
o & 2 & 1 & 1 & 2 & 3 \\
a & 3 & 2 & 1 & 2 & 3 \\
n & 4 & 3 & 2 & 1 & 2
\end{tabular}
\end{center}
The distance is 2: \code{loan-} over \code{l-ane}, deleting \code{o} and inserting \code{e}.
\end{solution}

\begin{solution}{exo:b1:dynamic-programming:2}
\code{fib_naive(6)} makes $2F(7) - 1 = 25$ calls, \code{fib_memo(6)} makes $2 \cdot 6 - 1 = 11$: one call for each $n$
from 6 down to 2 that computes, plus one call for each value read back or below 2.
\end{solution}

\begin{solution}{exo:b1:dynamic-programming:3}
From the last row $0, 8, 15, 25, 40, 50, 58, 65, 75$, the new entry for capacity $c$ is the larger of the old one and
30 plus the old one at $c - 3$: $0, 8, 15, 30, 40, 50, 58, 70, 80$. The best value is 80, with the items of weights 5 and
3 (values 50 and 30).
\end{solution}

\begin{solution}{exo:b1:dynamic-programming:4}
Tops after each term: $3$; $1$; $1, 4$; $1, 4$; $1, 4, 5$; $1, 4, 5, 9$; $1, 2, 5, 9$; $1, 2, 5, 6$; $1, 2, 5, 6$; $1, 2,
3, 6$; $1, 2, 3, 5$. Four piles: length 4, for instance $1, 4, 5, 9$ or $1, 2, 3, 5$.
\end{solution}

\begin{solution}{exo:b1:dynamic-programming:5}
Downwards, \code{best[c - w]} has not yet been updated for the current item: it is the best value without it, and each
item is used at most once, which gives 65. Upwards, \code{best[c - w]} may already contain the current item, which is
then taken again: 80 is twice the item $(4, 40)$. That is the right method for a knapsack with unlimited copies of each
item, the wrong one here.
\end{solution}

\begin{solution}{exo:b1:dynamic-programming:6}
The recursion always passes the same list object, shortened, so \code{id(values)} is the same at every depth: the result
for a prefix of three values is stored under the same key as the result for five, and read back in place of it. For
the total 10 the confusion happens to give the right count; for 13 and 17 it does not. The subproblem is a prefix and a
total, so the key should say which prefix: \code{(len(values), target)}, or \code{(tuple(values), target)}, which costs
a copy.
\end{solution}

\begin{solution}{exo:b1:dynamic-programming:7}
\omcode{code/bachelor-1/23-dynamic-programming/exercises/weighted_intervals.py}{6}{14}{Weighted
interval scheduling}
Sort the talks by finishing time; let $B(i)$ be the best weight using the first $i$, and $p(i)$ the number of talks among
the first $i - 1$ that finish by the start of talk $i$. An optimal set either leaves talk $i$ out, or takes it together
with an optimal set of the $p(i)$ talks compatible with it: $B(i) = \max(B(i - 1), B(p(i)) + w_i)$. Sorting and the
binary searches cost $O(n \log n)$. A heavy talk overlapping two light compatible ones defeats the earliest-finish rule.
\end{solution}

\begin{solution}{exo:b1:dynamic-programming:8}
\omcode{code/bachelor-1/23-dynamic-programming/exercises/stairs.py}{4}{8}{Ways to climb}
The last step of a climb to $k$ has some size $s$, and before it comes any climb to $k - s$: $W(k) = \sum_{s \leq k}
W(k - s)$ with $W(0) = 1$, in $O(n \cdot |S|)$. With steps 1 and 2, $W(n)$ is the Fibonacci number $F(n + 1)$: 89 for
$n = 10$.
\end{solution}

\begin{solution}{exo:b1:dynamic-programming:9}
An \omterm{def:b1:dynamic-programming:edit}{alignment} of cost $c$ is read as $c$ edits, column by column, so $d(a, b)$ is at most the smallest cost. Conversely,
take $d$ edits turning $a$ into $b$. Edits never reorder letters: the letters of $a$ that survive, perhaps changed, end
at positions of $b$ in the same order. Write each over its final letter, the deleted letters of $a$ over gaps and the
other letters of $b$ under gaps. A column of two different letters needs at least one substitution of that letter, a
column with a gap at least one deletion or insertion, all different edits: the \omterm{def:b1:dynamic-programming:edit}{alignment} costs at most $d$.
\end{solution}

\begin{solution}{exo:b1:dynamic-programming:10}
A key goes on a pile only if its top is not smaller, so each pile decreases from bottom to top; the keys of one pile came
in that order, so an increasing subsequence takes at most one key from each pile, and its length is at most the number
of piles. When a key $x$ goes on pile $k > 0$, the top $y$ of pile $k - 1$ is less than $x$ and came before it; remember
$y$. From any key on the last pile, following these remembered keys back to pile 0 gives keys that increase and come in
order, one per pile: an increasing subsequence as long as the number of piles.
\end{solution}

\begin{solution}{exo:b1:dynamic-programming:11}
\omcode{code/bachelor-1/23-dynamic-programming/exercises/palindrome.py}{4}{16}{The longest palindromic subsequence}
For $s[i..j]$: if the end letters are equal, a best palindrome may use both around a best one of $s[i + 1..j - 1]$;
otherwise it does without one of them. Cells are filled for $i$ decreasing and $j$ increasing, so each reads cells
already computed: $\Theta(n^2)$. The test compares with all subsequences of random words of up to 10 letters.
\end{solution}

\begin{solution}{exo:b1:dynamic-programming:12}
Without substitutions, every column of an \omterm{def:b1:dynamic-programming:edit}{alignment} of finite cost holds two equal letters, or a letter and a gap. The
columns of equal letters spell a common subsequence of some length $M$; the other $m - M$ letters of $a$ and $n - M$
of $b$ face gaps. So the cost is $m + n - 2M$, smallest when $M$ is the length $L$ of a \omterm{def:b1:dynamic-programming:lcs}{longest common subsequence}, and
by \cref{exo:b1:dynamic-programming:9} the distance is $m + n - 2L$.
\end{solution}

\begin{solution}{pb:b1:dynamic-programming:1}
\textbf{1.} One substitution, \code{i} into \code{e}, turns one into the other, and the words differ, so the distance is
not 0.

\textbf{2.} $20 \times 22 = 440$.

\textbf{3.} Row by row, each from left to right: a cell reads its left, upper and upper-left neighbours, all filled
before it.

\textbf{4.} $\Theta(mn)$, a constant amount of work per cell.

\textbf{5.} 80, one per title.

\textbf{6.} Turning $i$ letters into none takes $i$ deletions, and no operation removes more than one letter.

\textbf{7.} \code{the clo-kwork harbo-r} over \code{the clockwork harbour}: two insertions.

\textbf{8.} \code{catlaogue} over \code{catalogue} (two substitutions), and \code{cat-laogue} over \code{catal-ogue} (an
insertion and a deletion).

\textbf{9.} A cell's value is the smallest of three expressions, so it equals at least one of them.

\textbf{10.} $O(m + n)$: each step lowers $i + j$ by at least 1.

\textbf{11.} The walk back reads cells of every row, and only the last two rows are kept.

\textbf{12.} Compute the middle row from the top and the same row from the bottom, on the reversed words, each with two
rows; the column where their sum is smallest is where an optimal \omterm{def:b1:dynamic-programming:edit}{alignment} crosses the middle; recurse on the two
halves (Hirschberg's method).

\textbf{13.} 1258 cells against 68.

\textbf{14.} $1001^2 = 1\,002\,001$ against $2 \times 1001 = 2002$.

\textbf{15.} No: both versions compute every cell once, $\Theta(mn)$.

\textbf{16.} $4 \times 80 = 320$.

\textbf{17.} 45 against 640.

\textbf{18.} A query saves about $80 - 45/8 \approx 74$ distances, so the 320 pay off from the fifth query on.

\textbf{19.} No edit is needed exactly when the words are equal.

\textbf{20.} Each edit can be undone by one edit, an insertion by a deletion and a substitution by a substitution: the
reversed sequence turns $b$ into $a$ in as many steps.

\textbf{21.} Following $d(a, b)$ edits from $a$ to $b$, then $d(b, c)$ edits from $b$ to $c$, turns $a$ into $c$; the
smallest number of edits is at most that.

\textbf{22.} $d(q, p) \leq d(q, t) + d(t, p)$ gives $d(q, p) - d(p, t) \leq d(q, t)$, and $d(p, t) \leq d(p, q) + d(q, t)$
gives $d(p, t) - d(q, p) \leq d(q, t)$; for several pivots, the largest of these bounds.

\textbf{23.} The titles come by increasing bound, and each bound is at most the true distance: once a bound reaches the
best distance found, every remaining title is at least that far. (A title at an equal distance may be skipped: the
program returns one nearest title.)

\textbf{24.} \emph{\omterm{def:b1:dynamic-programming:edit}{Edit distance} is a metric}: it is zero exactly on equal words, symmetric, and satisfies the triangle
inequality (questions 19--21). On the catalogue, the bound it gives lets the eight queries compute 45 distances instead
of 640, after 320 computed once for the pivots.
\end{solution}
